\medskip

Un commerçant vend deux types de matelas: matelas RESSORTS et matelas MOUSSE. 

On suppose que chaque client achète un seul matelas.

On dispose des informations suivantes :
\begin{itemize}
\item[$\bullet~~$] 20\,\% des clients achètent un matelas RESSORTS. 

Parmi eux, 90\,\% sont satisfaits de leur achat.
\item[$\bullet~~$] 82\,\% des clients sont satisfaits de leur achat.
\end{itemize}



\medskip

\textbf{Partie A}

\medskip








On note $x = P_{\overline{R}}(S)$, où $P_{\overline{R}}(S)$ désigne la probabilité de $S$ sachant que $R$ n'est pas réalisé.

\begin{minipage}{0.58\linewidth}
\begin{enumerate}
\item Recopier et compléter l'arbre pondéré ci-contre décrivant la situation.
\end{enumerate}

\end{minipage}\hfill
\begin{minipage}{0.38\linewidth}
\begin{center}
\pstree[treemode=R,nodesepA=0pt,nodesepB=2.5pt,treesep = 1cm,levelsep=2.5cm]{\TR{}}
{\pstree{\TR{$R$~~}\taput{0,20}}
{\TR{$S$}\taput{0,9}
\TR{$\overline{S}$}\tbput{0,1}
}
\pstree{\TR{$\overline{R}$~~}\tbput{0,80}}
{\TR{$S$}\taput{$x$}
\TR{$\overline{S}$}\tbput{$1 - x$}
}}
\end{center}
\end{minipage}

\begin{enumerate}[start=2]
\item 
D'après la loi des probabilités totales :
$P(S) = P(R \cap S) + P\left(\overline{R} \cap S \right)$ Or :

$\bullet~~P(R \cap S) = P(R) \times P_R(S) = 0,2 \times 0,9 = 0,18$ et 

$\bullet~~P\left(\overline{R} \cap S \right) = P\left(\overline{R}\right) \times P_{\overline{R}}(S) = 0,8 \times x = 0,8x$ et comme $P(S) = 0,82$, on a donc l'équation :

$0,18 + 0,8x = 0,82 \iff 0,8x = 0,64 \iff x = 0,8$.
\item On choisit un client satisfait de son achat.

Quelle est la probabilité qu'il ait acheté un matelas RESSORTS ?


Il faut calculer $P_S(R) = \dfrac{P(S \cap R)}{P(S)} = \dfrac{0,18}{0,82} = \dfrac{18}{82} = \dfrac{9}{41} \approx 0,219$, soit 0,22 au centième près.
\end{enumerate}
\medskip

\medskip

\textbf{Partie B}

\medskip

\begin{enumerate}
\item On choisit 5 clients au hasard.

On considère la variable aléatoire $X$ qui donne le
nombre de clients satisfaits de leur achat parmi ces 5 clients.
	\begin{enumerate}
		\item 
Les paramètres de la loi binomiale sont $ n = 5 $ \quad et \quad  $ p=\np{0.82}$.
		\item 
		

Il faut trouver $P(X \leqslant 3) \approx \np{0,2223}$, d'après la calculatrice soit 0,222 au millième près.
	\end{enumerate}
\item 


	\begin{enumerate}
		\item On note $p_n$ la probabilité que les $n$ clients soient tous satisfaits de leur achat.


La probabilité qu'un acheteur soit satisfait est égale à 0,82, la probabilité que deux acheteurs soient satisfaits est $0,82^2$, \ldots, la probabilité que les $n$ acheteurs soient satisfaits est égale à $0,82^n$.
		\item 
		

Il faut donc résoudre dans $\N$ l'inéquation :

$p_n < 0,01 \iff 0,82^n < 0,01 \iff n \ln 0,82 < \ln 0,01$ par croissance de la fonction logarithme népérien soit encore 

$n > \dfrac{\ln 0,01}{\ln 0,82}$. (car $\ln 0,82 < 0$). Or $\dfrac{\ln 0,01}{\ln 0,82} \approx 23,2$.

Il faut donc au minimum 24 acheteurs.

Moins de 1\,\% des acheteurs seront tous satisfaits si leur nombre dépasse 23.
	\end{enumerate}
\end{enumerate}

\bigskip
